\chapter{Заключение}
\label{chap:conclusion}
\setlength{\parskip}{0.8em}

\section{Сводка вкладов}

В данной диссертации разработана, обучена и валидирована система автоматической детекции и классификации ожоговых поражений на основе архитектуры глубокого обучения YOLOv8. Используя датасет Kaggle Skin Burn Dataset ($n{=}1\,225$ изображений, 3 класса ожогов), модель достигла mAP@0.5~=~0,847, precision~=~0,823, recall~=~0,791.

Основные научные и практические вклады:

\textbf{1.~Методологический вклад.} Разработан воспроизводимый конвейер обработки медицинских изображений, включающий загрузку датасета из Kaggle, преобразование в YOLO-совместимый формат, аугментацию данных и обучение с трансферным обучением. Конвейер полностью автоматизирован и документирован.

\textbf{2.~Применение YOLOv8 для медицинской визуализации.} Впервые применена современная архитектура YOLOv8 для задачи детекции ожоговых поражений. В отличие от традиционных классификационных подходов (CNN), детекция обеспечивает одновременную классификацию и локализацию области поражения.

\textbf{3.~Практическая реализация.} Разработан REST API на базе фреймворка Litestar для интеграции модели в клинические информационные системы. API поддерживает обработку JPEG-изображений и возвращает структурированные результаты детекции в JSON-формате.

\textbf{4.~Инференс в реальном времени.} Выбранная архитектура YOLOv8n обеспечивает инференс со скоростью $>$20 FPS даже на CPU, что позволяет использовать модель в условиях ограниченных вычислительных ресурсов.

\textbf{5.~Облачное обучение.} Отработана методология обучения моделей глубокого обучения на облачной платформе vast.ai, обеспечивающая доступ к GPU-ресурсам для исследователей и студентов.

\section{Научная новизна: итоги}

Заявленные положения научной новизны (Глава~\ref{chap:intro}) были подтверждены:
\begin{itemize}
    \item Архитектура YOLOv8, изначально разработанная для детекции общих объектов, успешно адаптирована для специализированной медицинской задачи — классификации ожогов.
    \item Формула mAP (уравнение~\ref{eq:map}) применена для оценки качества детекции ожоговых поражений, обеспечивая объективное сравнение с другими подходами.
    \item Разработан комплексный конвейер аугментации данных, включающий HSV-преобразования, геометрические искажения и Mosaic-аугментацию, адаптированный для медицинских изображений.
\end{itemize}

\section{Практические рекомендации}

На основе результатов данной диссертации предлагаются следующие практические рекомендации:

\begin{enumerate}
    \item \textbf{Контекст применения.} Модель наилучшим образом подходит для предварительной оценки степени ожога как вспомогательный инструмент, дополняющий, но не заменяющий экспертную оценку врача-комбустиолога.

    \item \textbf{Телемедицина.} Система может использоваться для дистанционной консультации, когда пациент отправляет фотографию ожога, а врач получает предварительную автоматическую оценку степени.

    \item \textbf{Сортировка пациентов.} В условиях массового поступления пострадавших (ЧС, катастрофы) автоматическая классификация позволяет быстро определить приоритет лечения.

    \item \textbf{Обучение персонала.} Модель может использоваться для обучения медицинских сестёр и фельдшеров распознаванию степеней ожогов.
\end{enumerate}

\section{Направления будущих исследований}

Несколько перспективных расширений данной работы:

\begin{enumerate}
    \item \textbf{Расширение датасета.} Сбор дополнительных изображений, особенно ожогов третьей степени, для устранения дисбаланса классов. Рекомендуемый объём — не менее 5\,000 изображений.

    \item \textbf{Сегментация.} Переход от детекции (bounding box) к семантической сегментации позволит точнее определять площадь ожогового поражения, что критично для оценки TBSA (Total Body Surface Area).

    \item \textbf{Временная динамика.} Анализ серии фотографий одного ожога в разные моменты времени позволит прогнозировать прогрессирование или заживление.

    \item \textbf{Мультимодальность.} Интеграция дополнительных источников данных: термография, анамнез (причина ожога, время травмы), демографические данные пациента.

    \item \textbf{Мобильное приложение.} Разработка мобильного приложения для iOS/Android с встроенной моделью для offline-инференса на устройстве пациента.

    \item \textbf{Клиническая валидация.} Проспективное исследование точности модели в сравнении с экспертной оценкой комбустиологов на выборке казахстанских пациентов.

    \item \textbf{Интеграция с ИС.} Интеграция REST API с медицинскими информационными системами казахстанских больниц для автоматической документации ожоговых травм.
\end{enumerate}

\section{Заключительные замечания}

Применение глубокого обучения к классификации ожоговых поражений представляет научно обоснованный и практически значимый подход к совершенствованию медицинской помощи при ожоговых травмах. Данная диссертация демонстрирует, что современная архитектура детекции объектов YOLOv8, обученная на относительно небольшом специализированном датасете, может достичь клинически приемлемой точности классификации степеней ожогов.

Казахстанская система здравоохранения, активно внедряющая цифровые технологии, может существенно выиграть от интеграции подобных инструментов в экстренную и плановую медицинскую помощь. Разработанная методология обеспечивает переносимый шаблон для других задач медицинской визуализации: классификации кожных заболеваний, оценки заживления ран, диагностики дерматологических патологий.

Особую ценность представляет демонстрация возможности проведения исследований в области глубокого обучения с использованием облачных вычислительных ресурсов (vast.ai), что делает современные методы ИИ доступными для исследователей и студентов, не имеющих собственной GPU-инфраструктуры.

Конечным мерилом успеха этой работы станет её внедрение в клиническую практику — улучшение качества диагностики ожогов, сокращение времени принятия решений и, как результат, улучшение исходов лечения пациентов с ожоговыми травмами в Казахстане и за его пределами.
